% 附录B：傅里叶变换表

\chapter{傅里叶变换表}

\section{定义}
\begin{zhongyao}[傅里叶变换定义]
\begin{align}
\mathcal{F}[f(t)] &= \hat{f}(\omega) = \int_{-\infty}^{+\infty} f(t)\mathrm{e}^{-\mathrm{i}\omega t}\,\mathrm{d}t \\
\mathcal{F}^{-1}[\hat{f}(\omega)] &= f(t) = \frac{1}{2\pi}\int_{-\infty}^{+\infty} \hat{f}(\omega)\mathrm{e}^{\mathrm{i}\omega t}\,\mathrm{d}\omega
\end{align}
\end{zhongyao}

\section{常用傅里叶变换对}

\begin{table}[htbp]
\centering
\begin{tabular}{ll}
\toprule
$f(t)$ & $\hat{f}(\omega)$ \\
\midrule
$\mathrm{e}^{-a|t|}$ & $\dfrac{2a}{a^2+\omega^2}$ \\[8pt]
$\mathrm{e}^{-at^2}$ & $\sqrt{\dfrac{\pi}{a}}\mathrm{e}^{-\omega^2/(4a)}$ \\[8pt]
$\mathrm{e}^{-at}H(t)$ & $\dfrac{1}{a+\mathrm{i}\omega}$ \\[8pt]
$\delta(t)$ & $1$ \\[4pt]
$1$ & $2\pi\delta(\omega)$ \\[4pt]
$\cos\omega_0 t$ & $\pi[\delta(\omega-\omega_0) + \delta(\omega+\omega_0)]$ \\
$\sin\omega_0 t$ & $-\mathrm{i}\pi[\delta(\omega-\omega_0) - \delta(\omega+\omega_0)]$ \\
\bottomrule
\end{tabular}
\end{table}

\section{基本性质}

\begin{table}[htbp]
\centering
\begin{tabular}{ll}
\toprule
\textbf{性质} & \textbf{公式} \\
\midrule
线性性 & $\mathcal{F}[af + bg] = a\hat{f} + b\hat{g}$ \\
时移 & $\mathcal{F}[f(t-t_0)] = \mathrm{e}^{-\mathrm{i}\omega t_0}\hat{f}(\omega)$ \\
频移 & $\mathcal{F}[f(t)\mathrm{e}^{\mathrm{i}\omega_0 t}] = \hat{f}(\omega-\omega_0)$ \\
微分 & $\mathcal{F}[f'(t)] = \mathrm{i}\omega\hat{f}(\omega)$ \\
卷积 & $\mathcal{F}[f * g] = \hat{f}(\omega)\hat{g}(\omega)$ \\
帕塞瓦尔 & $\int |f|^2\,\mathrm{d}t = \frac{1}{2\pi}\int |\hat{f}|^2\,\mathrm{d}\omega$ \\
\bottomrule
\end{tabular}
\end{table}
